% GENERATED by md2tex.py from ../draft_v5_en/front/15_preface.md -- do not edit by hand.
% Change the Markdown and run `python md2tex.py`; edits here are overwritten.

This thesis is submitted as the Master's Thesis for the Automotive Systems (M.Eng.) programme at Hochschule Esslingen, supervised by Prof. Dr.-Ing. Ralf Schüler. The work starts from one idea: using artificial intelligence responsibly in automotive safety functions needs more than capable algorithms. It also needs an engineering method that makes their behaviour traceable, verifiable and auditable. The methodological framework developed in these pages, and applied to a concrete case, was built on that idea.

\textbf{Motivation and origin.} The project started from the question: if a learned component cannot be specified in advance and cannot be checked against an expected answer, what is left of the V-Model? The solution defended here is not to leave it as it is, but to modify it in such a way that the resulting cycle can continue to run from start to end.

\textbf{Scope.} The study goes from hazard analysis and requirement derivation to the validation campaign in simulation and the measurement of the gap towards the physical platform, on a single case study that is limited on purpose. The physical deployment is built, documented and running on hardware (the vehicle drives), but the physical campaign has not been carried out and no scenario has been scored on the real platform. The physical evidence in this document is bring-up evidence, not verdict evidence. Every run was made in monitoring mode, so the cage has never changed an action on the real vehicle, and for that reason the physical column of the final requirements table is marked not executed. Where there is no measurement, the document says so instead of estimating it. Where there is a measurement but it was not taken under protocol, it says that too.

\textbf{On honest reporting.} One decision runs through the whole work: uncomfortable results are reported as results. The global verdict of the reference campaign is literally \enquote{Not Satisfied}, one requirement is closed as not satisfied, and the conclusions chapter spends more space on what the framework did not show than on what it did. This is a choice of method: a traceability framework whose value is to leave no claim without evidence would lose its meaning if it were only applied to the favourable results.

I thank my supervisor for the guidance throughout the project, and Hochschule Esslingen for the resources it made available.
